%===============================================================================
% LaTeX sjabloon voor de bachelorproef toegepaste informatica aan HOGENT
% Meer info op https://github.com/HoGentTIN/latex-hogent-report
%===============================================================================

\documentclass[dutch,dit,thesis]{hogentreport}

% TODO:
% - If necessary, replace the option `dit`' with your own department!
%   Valid entries are dbo, dbt, dgz, dit, dlo, dog, dsa, soa
% - If you write your thesis in English (remark: only possible after getting
%   explicit approval!), remove the option "dutch," or replace with "english".

\usepackage{lipsum} % For blind text, can be removed after adding actual content

%% Pictures to include in the text can be put in the graphics/ folder
\graphicspath{{../graphics/}}

%% For source code highlighting, requires pygments to be installed
%% Compile with the -shell-escape flag!
%% \usepackage[chapter]{minted}
%% If you compile with the make_thesis.{bat,sh} script, use the following
%% import instead:
\usepackage{minted}
\usemintedstyle{solarized-light}

%% Formatting for minted environments.
\setminted{%
    autogobble,
    frame=lines,
    breaklines,
    linenos,
    tabsize=4
}

%% Ensure the list of listings is in the table of contents
\renewcommand\listoflistingscaption{%
    \IfLanguageName{dutch}{Lijst van codefragmenten}{List of listings}
}
\renewcommand\listingscaption{%
    \IfLanguageName{dutch}{Codefragment}{Listing}
}
\renewcommand*\listoflistings{%
    \cleardoublepage\phantomsection\addcontentsline{toc}{chapter}{\listoflistingscaption}%
    \listof{listing}{\listoflistingscaption}%
}

% Other packages not already included can be imported here

%%---------- Document metadata -------------------------------------------------
% TODO: Replace this with your own information
\author{Kristof Areias}
\supervisor{Jan Claes}
\cosupervisor{Thierry Klougbo}
\title[Optionele ondertitel]%
    {Automatische productherkenning in Vlaamse supermarkten met behulp van AI}
\academicyear{2025-2026}
\examperiod{1}
\degreesought{\IfLanguageName{dutch}{Professionele bachelor in de toegepaste informatica}{Bachelor of applied computer science}}
\partialthesis{false} %% To display 'in partial fulfilment'
%\institution{Internshipcompany BVBA.}

%% Add global exceptions to the hyphenation here
\hyphenation{back-slash}

%% The bibliography (style and settings are  found in hogentthesis.cls)
\addbibresource{bachproef.bib}            %% Bibliography file
\addbibresource{../voorstel/voorstel.bib} %% Bibliography research proposal
\defbibheading{bibempty}{}

%% Prevent empty pages for right-handed chapter starts in twoside mode
\renewcommand{\cleardoublepage}{\clearpage}

\renewcommand{\arraystretch}{1.2}

%% Content starts here.
\begin{document}

%---------- Front matter -------------------------------------------------------

\frontmatter

\hypersetup{pageanchor=false} %% Disable page numbering references
%% Render a Dutch outer title page if the main language is English
\IfLanguageName{english}{%
    %% If necessary, information can be changed here
    \degreesought{Professionele Bachelor toegepaste informatica}%
    \begin{otherlanguage}{dutch}%
       \maketitle%
    \end{otherlanguage}%
}{}

%% Generates title page content
\maketitle
\hypersetup{pageanchor=true}

%%=============================================================================
%% Voorwoord
%%=============================================================================

\chapter*{\IfLanguageName{dutch}{Woord vooraf}{Preface}}%
\label{ch:voorwoord}

%% TODO:
%% Het voorwoord is het enige deel van de bachelorproef waar je vanuit je
%% eigen standpunt (``ik-vorm'') mag schrijven. Je kan hier bv. motiveren
%% waarom jij het onderwerp wil bespreken.
%% Vergeet ook niet te bedanken wie je geholpen/gesteund/... heeft

In dit voorwoord wil ik mijn motivatie voor het gekozen onderwerp toelichten en mijn dank uitspreken aan de mensen die mij tijdens deze bachelorproef hebben ondersteund.
Toen ik dit onderwerp in de lijst zag staan, sprak het mij meteen aan. De combinatie van artificiële intelligentie en een praktische toepassing in de retailsector sluit goed aan bij mijn interesses. 
Vooral het idee om computer vision toe te passen in een realistische omgeving, zoals een supermarkt, maakte dit onderwerp voor mij extra boeiend.
Tijdens het uitwerken van deze bachelorproef heb ik niet alleen mijn technische kennis verder kunnen verdiepen, maar ook inzicht gekregen in de uitdagingen die komen kijken bij het toepassen van AI in de praktijk. Dit maakte het proces zowel leerrijk als uitdagend.
Daarnaast wil ik graag mijn promotor bedanken voor de begeleiding en feedback tijdens het traject. Ook wil ik mijn co-promotor bedanken voor zijn steun gedurende het hele proces.

%%=============================================================================
%% Samenvatting
%%=============================================================================

% TODO: De "abstract" of samenvatting is een kernachtige (~ 1 blz. voor een
% thesis) synthese van het document.
%
% Een goede abstract biedt een kernachtig antwoord op volgende vragen:
%
% 1. Waarover gaat de bachelorproef?
% 2. Waarom heb je er over geschreven?
% 3. Hoe heb je het onderzoek uitgevoerd?
% 4. Wat waren de resultaten? Wat blijkt uit je onderzoek?
% 5. Wat betekenen je resultaten? Wat is de relevantie voor het werkveld?
%
% Daarom bestaat een abstract uit volgende componenten:
%
% - inleiding + kaderen thema
% - probleemstelling
% - (centrale) onderzoeksvraag
% - onderzoeksdoelstelling
% - methodologie
% - resultaten (beperk tot de belangrijkste, relevant voor de onderzoeksvraag)
% - conclusies, aanbevelingen, beperkingen
%
% LET OP! Een samenvatting is GEEN voorwoord!

%%---------- Nederlandse samenvatting -----------------------------------------
%
% TODO: Als je je bachelorproef in het Engels schrijft, moet je eerst een
% Nederlandse samenvatting invoegen. Haal daarvoor onderstaande code uit
% commentaar.
% Wie zijn bachelorproef in het Nederlands schrijft, kan dit negeren, de inhoud
% wordt niet in het document ingevoegd.

% \IfLanguageName{english}{%
% \selectlanguage{dutch}
% \chapter*{Samenvatting}
% \lipsum[1-4]
% \selectlanguage{english}
% }{}

%%---------- Samenvatting -----------------------------------------------------
% De samenvatting in de hoofdtaal van het document

\chapter*{\IfLanguageName{dutch}{Samenvatting}{Abstract}}

De transparantie over het voedingsaanbod in Vlaamse supermarkten is beperkt, hoewel gezonde en duurzame voeding steeds centraler staan in beleid en onderzoek. 
Er bestaat geen geautomatiseerde manier om het werkelijke assortiment in kaart te brengen, waardoor onderzoekers, beleidsmakers en consumenten onvoldoende zicht hebben op de aanwezigheid van productcategorieën, duurzame alternatieven en marktverhoudingen. 
Dit onderzoek evalueert bestaande computer vision-modellen (bijv. YOLO-achtige object detectors, multimodale OCR- en beeld-classificatiemodellen) op de taak van automatische productherkenning in beeldmateriaal uit Vlaamse supermarkten, met als doel de nauwkeurigheid van detectie te beoordelen en eventueel te verbeteren via fine-tuning.
Het doel is een reproduceerbare methodologie te ontwikkelen die beelddata verwerkt en productinformatie herkent. 
De aanpak omvat het verzamelen van representatieve beelddata, het testen van modellen zoals YOLOv8 en Grounding DINO, het analyseren van prestaties via mAP en IoU-gebaseerde detectiemetrics, aangevuld met SKU-niveau Top-1 en Top-5 classificatie-accuratesse, macro-F1-scores en inventarisatiegerichte indicatoren zoals product coverage rate en assortimentsafwijking. 
Verwacht wordt dat deze studie de haalbaarheid aantoont van automatische inventarisatie in een complex retailecosysteem en dat de resultaten bijdragen aan betere monitoring van duurzaamheid en gezondheid, met concrete meerwaarde voor onderzoekers, beleidsmakers, consumenten en retailers.



%---------- Inhoud, lijst figuren, ... -----------------------------------------

\tableofcontents

% In a list of figures, the complete caption will be included. To prevent this,
% ALWAYS add a short description in the caption!
%
%  \caption[short description]{elaborate description}
%
% If you do, only the short description will be used in the list of figures

\listoffigures

% If you included tables and/or source code listings, uncomment the appropriate
% lines.
\listoftables

\listoflistings

% Als je een lijst van afkortingen of termen wil toevoegen, dan hoort die
% hier thuis. Gebruik bijvoorbeeld de ``glossaries'' package.
% https://www.overleaf.com/learn/latex/Glossaries

%---------- Kern ---------------------------------------------------------------

\mainmatter{}

% De eerste hoofdstukken van een bachelorproef zijn meestal een inleiding op
% het onderwerp, literatuurstudie en verantwoording methodologie.
% Aarzel niet om een meer beschrijvende titel aan deze hoofdstukken te geven of
% om bijvoorbeeld de inleiding en/of stand van zaken over meerdere hoofdstukken
% te verspreiden!

%%=============================================================================
%% Inleiding
%%=============================================================================

\chapter{\IfLanguageName{dutch}{Inleiding}{Introduction}}%
\label{ch:inleiding}

De samenstelling van het voedingsaanbod in supermarkten speelt een cruciale rol in consumentengedrag, gezondheid en duurzaamheid. Beleidsinitiatieven rond gezonde voeding en ecologische impact vertrekken vaak vanuit productinformatie zoals voedingswaarden, labels of duurzaamheidsindicatoren. Toch ontbreekt in Vlaanderen een systematische en schaalbare methode om het werkelijke schapaanbod objectief te registreren.

Traditionele inventarisatiemethoden zijn voornamelijk manueel: onderzoekers noteren producten ter plaatse of verwerken kassadata indien beschikbaar. Deze aanpak is arbeidsintensief, beperkt in schaal en moeilijk reproduceerbaar. Bovendien geven kassadata enkel inzicht in verkochte producten, niet in het volledige aangeboden assortiment.

Tegelijkertijd hebben recente ontwikkelingen binnen computer vision aangetoond dat objectdetectie- en beeldclassificatiemodellen in staat zijn om objecten in complexe visuele omgevingen te lokaliseren en te identificeren. Architecturen zoals YOLO-gebaseerde detectiemodellen en transformer-gebaseerde modellen tonen sterke prestaties in generieke beeldherkenningstaken. De toepassing van dergelijke modellen in realistische supermarktcontexten blijft echter uitdagend door factoren zoals hoge rekdichtheid, overlappende producten, variabele lichtomstandigheden en sterk gelijkende verpakkingen.

Deze bachelorproef onderzoekt of bestaande AI-modellen voldoende performant en robuust zijn om ingezet te worden voor automatische productdetectie in Vlaamse supermarktbeelden. De focus ligt expliciet op technische haalbaarheid en reproduceerbare evaluatie, niet op commerciële implementatie of beleidsanalyse. Door modellen systematisch te vergelijken onder gelijke omstandigheden wordt nagegaan of automatische assortimentsinventarisatie via beelddata een realistische onderzoeksstrategie kan vormen.

\section{\IfLanguageName{dutch}{Probleemstelling}{Problem Statement}}%
\label{sec:probleemstelling}

Onderzoekers binnen voedings- en duurzaamheidsstudies beschikken momenteel niet over een schaalbare methode om het feitelijke supermarktassortiment systematisch in kaart te brengen. Bestaande studies maken gebruik van manuele inventarisaties of steekproeven, wat tijdsintensief, foutgevoelig en moeilijk reproduceerbaar is. 

Voor Vlaamse onderzoeksgroepen die het aanbod van gezonde of duurzame producten willen monitoren (bijvoorbeeld in het kader van Nutri-Score, eco-score of productdiversiteit), ontbreekt een geautomatiseerde beeldgebaseerde aanpak die toelaat om op regelmatige basis het schapaanbod te analyseren.

Het gebrek aan een dergelijke methode beperkt:
\begin{itemize}
    \item zien hoe het aanbod in de winkel evolueert,
    \item objectieve vergelijking tussen supermarktlocaties,
    \item schaalbare analyse van productcategorieën.
\end{itemize}

Deze bachelorproef richt zich specifiek op de technische haalbaarheid van automatische productherkenning als eerste stap richting geautomatiseerde assortimentsanalyse.

\section{\IfLanguageName{dutch}{Onderzoeksvraag}{Research question}}%
\label{sec:onderzoeksvraag}

De centrale onderzoeksvraag luidt:

\begin{quote}
\textbf{In welke mate kunnen bestaande AI-modellen, al dan niet met bijkomende fine-tuning, producten in Vlaamse supermarkt, beelden correct herkennen als eerste stap richting geautomatiseerde assortimentsinventarisatie?}
\end{quote}

Om deze onderzoeksvraag te beantwoorden, wordt het probleem verder opgesplitst in een aantal deelvragen binnen het probleemdomein:

\begin{itemize}
    \item Wat is het effect van gecontroleerde variaties in belichting (lux-niveau), beeldresolutie (px) en rekdichtheid (aantal producten per m²) op de mAP- en recall-scores van geselecteerde detectiemodellen?
    \item Welke bestaande open datasets of online databronnen bevatten geschikt beeldmateriaal van supermarktproducten voor training en evaluatie van productherkenningsmodellen?
\end{itemize}

Daarnaast worden deelvragen geformuleerd binnen het oplossingsdomein, die focussen op de technische haalbaarheid van automatische productherkenning:

\begin{itemize}
    \item Welke mAP@0.5:0.95-, recall en precisionscores behalen geselecteerde AI-modellen bij detectie van producten in supermarktbeelden met hoge rekdichtheid?
    \item Hoe verschillen YOLO-modellen in termen van mAP, recall en detectiesnelheid bij identieke supermarktbeelden?
    \item Wat is het effect van beperkte fine-tuning op mAP, recal en precision ten opzichte van het vooraf getrainde basismodel?
\end{itemize}


\section{\IfLanguageName{dutch}{Onderzoeksdoelstelling}{Research objective}}%
\label{sec:onderzoeksdoelstelling}

De doelstelling van deze bachelorproef is het uitvoeren van een technisch onderbouwde vergelijkende studie naar de prestaties van bestaande computer vision-modellen voor productherkenning in Vlaamse supermarktbeelden.

Het beoogde resultaat is:
\begin{itemize}
    % \item een reproduceerbare experimentele evaluatie-opzet,
    \item een kwantitatieve vergelijking van modelprestaties,
    \item een proof-of-concept pipeline die detectie en metadata-koppeling integreert,
    \item concrete aanbevelingen over minimale prestatiedrempels voor inzetbaarheid in assortimentsmonitoring.
\end{itemize}

Het onderzoek wordt als succesvol beschouwd wanneer:
\begin{itemize}
    \item modelprestaties reproduceerbaar gemeten kunnen worden,
    \item verschillen tussen modellen statistisch aantoonbaar zijn,
    \item we objectief kunnen nagaan of de methode bruikbaar is voor het inventariseren van assortimenten.
\end{itemize}

\section{\IfLanguageName{dutch}{Opzet van deze bachelorproef}{Structure of this bachelor thesis}}%
\label{sec:opzet-bachelorproef}

% Het is gebruikelijk aan het einde van de inleiding een overzicht te
% geven van de opbouw van de rest van de tekst. Deze sectie bevat al een aanzet
% die je kan aanvullen/aanpassen in functie van je eigen tekst.

De rest van deze bachelorproef is als volgt opgebouwd:

In Hoofdstuk~\ref{ch:stand-van-zaken} bespreekt de stand van zaken binnen automatische productherkenning en objectdetectie in retailomgevingen. Hier worden relevante modelarchitecturen, evaluatiemetrieken en gekende uitdagingen in supermarktcontexten toegelicht.

In Hoofdstuk~\ref{ch:methodologie} beschrijft de experimentele opzet van het onderzoek. Dit omvat de selectie en voorbereiding van beelddata, de configuratie van de gekozen AI-modellen, de operationalisering van evaluatiemetrieken en de experimentele vergelijkingsstrategie.

% TODO: Vul hier aan voor je eigen hoofstukken, één of twee zinnen per hoofdstuk

In Hoofdstuk~\ref{ch:conclusie}, tenslotte, wordt de conclusie gegeven en een antwoord geformuleerd op de onderzoeksvragen. Daarbij wordt ook een aanzet gegeven voor toekomstig onderzoek binnen dit domein.
\chapter{\IfLanguageName{dutch}{Stand van zaken}{State of the art}}%
\label{ch:stand-van-zaken}

% Tip: Begin elk hoofdstuk met een paragraaf inleiding die beschrijft hoe
% dit hoofdstuk past binnen het geheel van de bachelorproef. Geef in het
% bijzonder aan wat de link is met het vorige en volgende hoofdstuk.

% Pas na deze inleidende paragraaf komt de eerste sectiehoofding.

% Dit hoofdstuk bevat je literatuurstudie. De inhoud gaat verder op de inleiding, maar zal het onderwerp van de bachelorproef *diepgaand* uitspitten. De bedoeling is dat de lezer na lezing van dit hoofdstuk helemaal op de hoogte is van de huidige stand van zaken (state-of-the-art) in het onderzoeksdomein. Iemand die niet vertrouwd is met het onderwerp, weet nu voldoende om de rest van het verhaal te kunnen volgen, zonder dat die er nog andere informatie moet over opzoeken \autocite{Pollefliet2011}.

% Je verwijst bij elke bewering die je doet, vakterm die je introduceert, enz.\ naar je bronnen. In \LaTeX{} kan dat met het commando \texttt{$\backslash${textcite\{\}}} of \texttt{$\backslash${autocite\{\}}}. Als argument van het commando geef je de ``sleutel'' van een ``record'' in een bibliografische databank in het Bib\LaTeX{}-formaat (een tekstbestand). Als je expliciet naar de auteur verwijst in de zin (narratieve referentie), gebruik je \texttt{$\backslash${}textcite\{\}}. Soms is de auteursnaam niet expliciet een onderdeel van de zin, dan gebruik je \texttt{$\backslash${}autocite\{\}} (referentie tussen haakjes). Dit gebruik je bv.~bij een citaat, of om in het bijschrift van een overgenomen afbeelding, broncode, tabel, enz. te verwijzen naar de bron. In de volgende paragraaf een voorbeeld van elk.

% \textcite{Knuth1998} schreef een van de standaardwerken over sorteer- en zoekalgoritmen. Experten zijn het erover eens dat cloud computing een interessante opportuniteit vormen, zowel voor gebruikers als voor dienstverleners op vlak van informatietechnologie~\autocite{Creeger2009}.

% Let er ook op: het \texttt{cite}-commando voor de punt, dus binnen de zin. Je verwijst meteen naar een bron in de eerste zin die erop gebaseerd is, dus niet pas op het einde van een paragraaf.

% \begin{figure}
%   \centering
%   \includegraphics[width=0.8\textwidth]{grail.jpg}
%   \caption[Voorbeeld figuur.]{\label{fig:grail}Voorbeeld van invoegen van een figuur. Zorg altijd voor een uitgebreid bijschrift dat de figuur volledig beschrijft zonder in de tekst te moeten gaan zoeken. Vergeet ook je bronvermelding niet!}
% \end{figure}

% \begin{listing}
%   \begin{minted}{python}
%     import pandas as pd
%     import seaborn as sns

%     penguins = sns.load_dataset('penguins')
%     sns.relplot(data=penguins, x="flipper_length_mm", y="bill_length_mm", hue="species")
%   \end{minted}
%   \caption[Voorbeeld codefragment]{Voorbeeld van het invoegen van een codefragment.}
% \end{listing}

% \lipsum[7-20]

% \begin{table}
%   \centering
%   \begin{tabular}{lcr}
%     \toprule
%     \textbf{Kolom 1} & \textbf{Kolom 2} & \textbf{Kolom 3} \\
%     $\alpha$         & $\beta$          & $\gamma$         \\
%     \midrule
%     A                & 10.230           & a                \\
%     B                & 45.678           & b                \\
%     C                & 99.987           & c                \\
%     \bottomrule
%   \end{tabular}
%   \caption[Voorbeeld tabel]{\label{tab:example}Voorbeeld van een tabel.}
% \end{table}

Dit hoofdstuk bouwt voort op de inleiding en plaatst het onderzoek in de bredere wetenschappelijke context van automatische productherkenning in retailomgevingen. Waar het vorige hoofdstuk het probleem en de onderzoeksvraag definieerde, bespreekt dit hoofdstuk bestaande datasets, methodes en manieren om resultaten te evalueren binnen grocery product recognition. Deze literatuurstudie vormt de theoretische basis voor de experimentele opzet die in het volgende hoofdstuk wordt beschreven.

\subsection*{1. Productherkenning als detectie- en identificatieprobleem}

Automatische productherkenning in supermarktcontext kan worden opgevat als een tweestapsprobleem bestaande uit (1) lokalisatie van producten in een schapbeeld en (2) identificatie op merk- of SKU-niveau. Volgens \textcite{MELEK2024108452} moet een goed werkend systeem niet alleen objecten correct detecteren, maar ze ook tot op productniveau kunnen onderscheiden (fine-grained classificatie). Dit is nodig omdat toepassingen in de retail vaak een hoge nauwkeurigheid per individueel product vereisen.

In tegenstelling tot algemene objectdetectietaken (zoals COCO), hebben supermarktbeelden enkele specifieke kenmerken:

\begin{itemize}
\item er staan vaak veel producten dicht bij elkaar,
\item producten lijken sterk op elkaar,
\item producten overlappen elkaar regelmatig,
\item verschillen tussen varianten van hetzelfde product zijn vaak klein.
\end{itemize}

\textcite{FRANCO2017163} tonen aan dat een belangrijke aanname uit klassieke objectdetectie vaak niet klopt in de retailcontext. Meestal gaat men ervan uit dat trainingsdata goed overeenkomt met de situatie in de praktijk. In supermarkten is dat echter vaak niet het geval. Trainingsbeelden worden genomen in gecontroleerde omstandigheden, terwijl beelden in winkels veel variatie en beweging bevatten. Daardoor ontstaat er een verschil tussen training en gebruik in de praktijk.

Productherkenning moet daarom niet alleen goed omgaan met overlappende producten en verschillen in grootte, maar ook met producten die sterk op elkaar lijken.

\subsection*{2. Datasetkarakteristieken en representativiteit}

Datasetkwaliteit vormt een fundamentele beperkende factor binnen grocery product recognition. De systematische review van \textcite{MELEK2024108452} toont aan dat bestaande datasets op meerdere punten tekortschieten:

\begin{itemize}
\item er is weinig variatie in het aantal producten,
\item beelden worden vaak genomen onder gecontroleerde belichting en vanuit vaste camerahoeken,
\item er staan relatief weinig producten in één beeld,
\item het annoteren van producten met bounding boxes is tijdsintensief en duur.
\end{itemize}

Daardoor lijken resultaten op deze datasets vaak beter dan ze in een echte winkelomgeving zouden zijn.

\textcite{TONIONI201810733} tonen ook aan dat herkenning op basis van embeddings sterk afhangt van de trainingsdata. Als variaties zoals belichting, camerahoek of productrotatie onvoldoende aanwezig zijn, daalt de performantie merkbaar.

Daarnaast tonen recente studies aan dat domeinadaptatie belangrijk is om het verschil tussen gecontroleerde datasets en echte winkelbeelden te verkleinen \autocite{wei2020retail}. Toch zijn er weinig grote en actuele retaildatasets beschikbaar. Dit komt onder andere door commerciële beperkingen en privacyredenen.

De kwaliteit en representativiteit van datasets is dus geen detail, maar heeft een grote invloed op hoe goed modellen presteren.

\subsection*{3. Klassieke computer vision benaderingen}

Voor de opkomst van deep learning maakten productherkenningssystemen gebruik van handgemaakte visuele kenmerken zoals:
\begin{itemize}
\item SURF-descriptoren; voor opvallende details zoals hoeken en randen,
\item kleurhistogrammen; voor de kleurverdeling van producten zoals welke kleuren er aanwezig zijn en in welke verhouding,
\item textuurfeatures; voor patronen in verpakkingen zoals verpakkingstructuren,
\item Bag-of-Visual-Words representaties; voor een visuele representatie zoals een verzmaling van de belangrijkste visuele kenmerken.
\end{itemize}

\textcite{FRANCO2017163} vergelijken deze klassieke technieken met deep learning features in een systeem met meerdere stappen, zoals pre-selectie, fine-selectie en post-processing. Hun resultaten tonen aan dat kleur- en textuurinformatie belangrijk blijven in een retailcontext. Tegelijk blijkt dat klassieke methoden moeilijk opschalen wanneer het aantal producten sterk toeneemt.

Hoewel deze technieken historisch belangrijk zijn, werken ze minder goed bij grote productassortimenten en wanneer producten binnen dezelfde categorie sterk van elkaar verschillen.

\subsection*{4. CNN-gebaseerde objectdetectie en embedding-methoden}

Met de introductie van convolutionele neurale netwerken (CNN’s) verschoof de focus naar modellen die end-to-end getraind kunnen worden.

Volgens \textcite{MELEK2024108452} domineren CNN-gebaseerde detectiearchitecturen het huidige retailonderzoek. Deze modellen leren zelf kenmerken uit de data en gaan beter om met verschillen in schaal en rotatie dan klassieke methoden.

\textcite{TONIONI201810733} stellen een aanpak voor waarbij detectie en herkenning apart gebeuren. Eerst worden producten gelokaliseerd in het beeld, daarna wordt via embedding-gebaseerde matching bepaald om welk product het precies gaat. Hierdoor wordt het probleem herleid tot het zoeken naar het meest gelijkaardige product in een feature-ruimte.

Hoewel CNN-gebaseerde systemen duidelijke prestatievoordelen bieden, blijven ze gevoelig voor:

\begin{itemize}
\item kleine producten,
\item overlappende producten,
\item visueel sterk gelijkende verpakkingen,
\item veranderende verpakkingdesigns.
\end{itemize}

Ook hier blijft de kwaliteit van de trainingsdata een belangrijke factor.

\subsection*{5. Transformer-gebaseerde detectiearchitecturen}

Recente ontwikkelingen binnen objectdetectie introduceren transformer-gebaseerde modellen als alternatief voor traditionele anchor-based detectie.

Het Detection Transformer (DETR)-model herformuleert detectie als een set-predictietaak zonder non-maximum suppression \autocite{carion2020}. Door bipartite matching wordt een één-op-één correspondentie geleerd tussen voorspellingen en ground-truth objecten.

Deformable DETR optimaliseert deze aanpak door efficiëntere attention-mechanismen te introduceren, wat leidt tot snellere convergentie en betere prestaties bij kleine objecten \autocite{zhu2021}. Dit is bijzonder relevant voor supermarktbeelden met hoge objectdensiteit.

Transformer-gebaseerde modellen bieden theoretische voordelen in complexe scenestructuren, maar vereisen doorgaans grotere trainingsdatasets.

\subsection*{6. Fine-grained classificatie op SKU-niveau}

Productherkenning in retail gaat verder dan gewone objectdetectie, omdat producten tot op SKU-niveau herkend moeten worden. Dit betekent dat het systeem zeer kleine visuele verschillen moet kunnen onderscheiden.

\textcite{PETTERSSON202410549} tonen aan dat visuele embeddings vaak niet voldoende onderscheid maken wanneer verpakkingen sterk op elkaar lijken. Kleine verschillen, zoals tekst, volume-aanduidingen of kleuraccenten, kunnen nochtans belangrijk zijn om producten correct te identificeren.

Daarom is het nodig om representaties te gebruiken die gevoelig zijn voor subtiele verschillen, zonder dat ze te sterk reageren op irrelevante details.

\subsection*{7. Multimodale benaderingen met beeld en tekst}

Om verwarring tussen sterk gelijkende producten te verminderen, kan men beeld- en tekstinformatie combineren.

Volgens \textcite{PETTERSSON202410549} zorgt het toevoegen van OCR-teksten aan visuele kenmerken voor een duidelijke verbetering van SKU-herkenning. Tekst zoals merk en producttype helpt om producten correct te identificeren.

Vision-language modellen zoals CLIP leren een gezamenlijke ruimte voor beeld en tekst. Daardoor kunnen tekstbeschrijvingen eenvoudig vergeleken worden met beeldfeatures.

Dergelijke multimodale modellen zijn robuuster tegen overlappende producten en visuele verwarring.

% \subsection*{8. Open-vocabulary en zero-shot detectie}

% Traditionele detectiemodellen werken vaak met een vaste set categorieën. In retail verandert het assortiment echter voortdurend, omdat er regelmatig nieuwe producten worden geïntroduceerd.

% Open-vocabulary detectie combineert visuele detectie met tekstuele aanwijzingen. Met Grounding DINO kunnen objecten worden gedetecteerd op basis van vrije tekstbeschrijvingen \autocite{wei2020retail}, waardoor het model minder afhankelijk is van vaste labels.

% Zero-shot herkenning via vision-language modellen maakt het mogelijk om producten van niet-getrainde categorieën te classificeren \autocite{radford2021clip}. Dit opent mogelijkheden voor monitoring van dynamische assortimenten.

% De uitdagingen van open-set herkenning worden besproken in de literatuur \autocite{peng2021rp2klargescaleretailproduct}, waar wordt benadrukt dat modellen onbekende klassen moeten kunnen detecteren zonder verkeerde toewijzingen.

\subsection*{9. Generaliseerbaarheid en domeinadaptatie}

In retailherkenning is domeinverschuiving een terugkerend probleem.

Volgens \textcite{FRANCO2017163} presteren modellen minder goed als testbeelden afwijken in belichting of camerastandpunt. Data-augmentatie en domeinadversarial training kunnen helpen om modellen beter te laten generaliseren \autocite{wei2020retail}, maar ze lossen het probleem niet helemaal op.

Het is daarom belangrijk om de generaliseerbaarheid te evalueren onder verschillende beeldomstandigheden, niet alleen op uniforme testsets.

\subsection*{10. Evaluatiekaders voor retaildetectie}

Evaluatie van productherkenning vereist een combinatie van detectie- en classificatiemetrieken.

Mean Average Precision (mAP) wordt gebruikt om lokalisatiekwaliteit te meten \autocite{MELEK2024108452}. mAP kijkt niet alleen naar of een product is gevonden, maar ook of het correct is omkaderd en hoe betrouwbaar de voorspellingen zijn over meerdere drempels. Andere metrics zoals IoU of F1-score per object zijn te beperkt: ze geven alleen informatie per detectie en samenvatten prestaties niet over categorieën of drempels. Daarom is mAP de standaard in situaties waar het gaat om het lokaliseren van vele SKU’s in beelden.

Recall meet hoeveel producten daadwerkelijk worden gevonden. In retailtoepassingen is dit cruciaal, omdat het missen van een product (false negative) vaak grotere gevolgen heeft dan een fout-positief detectie, bijvoorbeeld bij voorraadmonitoring of prijscontrole \autocite{FRANCO2017163}. Precision wordt minder zwaar meegewogen, omdat enkele foutieve detecties relatief makkelijk handmatig of automatisch gecorrigeerd kunnen worden. Andere metrics zoals F1-score wegen false positives en false negatives gelijk, maar dat weerspiegelt niet de praktijkbehoefte in winkels.

% Voor fine-grained classificatie worden Top-1 en Top-5 accuracy gebruikt \autocite{PETTERSSON202410549}. Top-1 laat zien of het model exact het juiste product (SKU) voorspelt. Top-5 is nuttig wanneer producten visueel sterk op elkaar lijken; het geeft aan of het juiste product in de vijf beste voorspellingen zit. Andere metrics zoals macro-F1 of gemiddelde accuracy over alle categorieën zijn minder informatief bij grote hoeveelheden SKU’s, omdat zij de prestatie per klasse gelijk wegen en de impact van veelvoorkomende producten onderschatten.

Bij open-set scenario’s moet gekeken worden naar foutieve toewijzing van onbekende producten \autocite{peng2021rp2klargescaleretailproduct}. Klassieke metrics zoals mAP negeren onbekende klassen volledig, waardoor ze onvoldoende weergeven hoe goed een model onbekende producten kan herkennen zonder deze verkeerd in te delen. Dit is van groot belang in dynamische retailomgevingen, waar regelmatig nieuwe producten verschijnen.

\subsection*{11. Geïdentificeerde tekortkomingen in de literatuur}

Uit de literatuur blijkt dat er verschillende belangrijke hiaten bestaan in onderzoek naar automatische productherkenning in supermarkten.

Ten eerste werken veel studies met gecontroleerde datasets die niet lijken op echte winkelomstandigheden, waar producten dicht op elkaar staan, belichting varieert en verpakkingen overlappen \autocite{MELEK2024108452}.

Ten tweede worden verschillende modellen zelden onder identieke omstandigheden vergeleken. De meeste studies testen één model of pipeline \autocite{TONIONI201810733}, waardoor we niet goed kunnen zien welk type model echt beter presteert.

Daarnaast is er weinig onderzoek naar de invloed van beeldfactoren zoals camerahoek, resolutie, afstand tot het rek of productdichtheid. Domeinverschuiving wordt wel genoemd \autocite{FRANCO2017163}, maar vaak ontbreekt kwantitatief bewijs over welke factoren detectieprestaties het meest beïnvloeden.

Ook blijft de evaluatie van nieuwe of onbekende producten beperkt. Modellen worden vaak getest op een vaste set categorieën \autocite{peng2021rp2klargescaleretailproduct}, terwijl assortimenten in supermarkten voortdurend veranderen.

Kortom, ondanks vooruitgang in technologie is er nog steeds nood aan systematische, reproduceerbare evaluaties in echte winkelomstandigheden.

\subsection*{12. Positionering van dit onderzoek binnen de literatuur}

Deze bachelorproef onderzoekt hoe goed bestaande AI-modellen werken voor het automatisch detecteren van producten in echte supermarktbeelden.

In tegenstelling tot studies die een nieuwe architectuur of pipeline voorstellen, ligt de focus hier op:

\begin{itemize}
\item de experimentele vergelijking van verschillende modeltypes (klassieke objectdetectie, transformer-gebaseerde detectie en multimodale benaderingen),
\item de analyse van detectieprestaties onder variërende beeldcondities,
\item de evaluatie van generaliseerbaarheid binnen een concrete Vlaamse supermarktcontext.
\end{itemize}

Door verschillende modellen onder dezelfde experimentele omstandigheden te testen, geeft dit onderzoek een beter inzicht in hun sterktes en beperkingen in drukke retailomgevingen.

Daarmee sluit het aan bij de nood aan contextspecifieke en reproduceerbare evaluaties van detectiemodellen, zoals beschreven in de literatuur \autocite{MELEK2024108452}.
%%=============================================================================
%% Methodologie
%%=============================================================================

\chapter{\IfLanguageName{dutch}{Methodologie}{Methodology}}%
\label{ch:methodologie}

%% TODO: In dit hoofstuk geef je een korte toelichting over hoe je te werk bent
%% gegaan. Verdeel je onderzoek in grote fasen, en licht in elke fase toe wat
%% de doelstelling was, welke deliverables daar uit gekomen zijn, en welke
%% onderzoeksmethoden je daarbij toegepast hebt. Verantwoord waarom je
%% op deze manier te werk gegaan bent.
%% 
%% Voorbeelden van zulke fasen zijn: literatuurstudie, opstellen van een
%% requirements-analyse, opstellen long-list (bij vergelijkende studie),
%% selectie van geschikte tools (bij vergelijkende studie, "short-list"),
%% opzetten testopstelling/PoC, uitvoeren testen en verzamelen
%% van resultaten, analyse van resultaten, ...
%%
%% !!!!! LET OP !!!!!
%%
%% Het is uitdrukkelijk NIET de bedoeling dat je het grootste deel van de corpus
%% van je bachelorproef in dit hoofstuk verwerkt! Dit hoofdstuk is eerder een
%% kort overzicht van je plan van aanpak.
%%
%% Maak voor elke fase (behalve het literatuuronderzoek) een NIEUW HOOFDSTUK aan
%% en geef het een gepaste titel.

Dit hoofdstuk beschrijft de algemene onderzoeksaanpak die werd gevolgd om de onderzoeksvragen te beantwoorden. Het onderzoek werd gefaseerd uitgevoerd, waarbij elke fase een duidelijk afgebakende doelstelling en concrete output had. Deze studie combineert literatuuronderzoek met experimentele tests in een realistische retailomgeving.
\section*{Onderzoeksopzet}

Het onderzoek werd opgebouwd volgens zes opeenvolgende fasen:

\begin{enumerate}
    \item Literatuurstudie
    \item Requirementsanalyse
    \item Selectie van modellen en datasets
    \item Ontwikkeling van een proof-of-concept (PoC)
    \item Experimentele evaluatie
    \item Analyse en conclusie
\end{enumerate}

Deze gefaseerde aanpak werd gekozen om eerst een theoretische onderbouwing te verzekeren alvorens over te gaan tot experimentele implementatie en evaluatie.

\section*{Fase 1: Literatuurstudie}

De eerste fase bestond uit een diepgaande literatuurstudie naar automatische productherkenning in supermarktcontext. Het doel van deze fase was:

\begin{itemize}
    \item inzicht verwerven in bestaande detectie- en herkenningsmethoden,
    \item identificeren van relevante datasets,
    \item analyseren van evaluatiemetrieken,
    \item detecteren van gebrekken in het huidige onderzoekslandschap.
\end{itemize}

De literatuurstudie vormde de theoretische basis voor de verdere methodologische keuzes binnen dit onderzoek.

\section*{Fase 2: Requirementsanalyse}

Op basis van de literatuur werd een requirementsanalyse opgesteld. Het doel hiervan was het expliciet definiëren van:

\begin{itemize}
    \item functionele vereisten (bv. correcte detectie van producten in dense schapbeelden),
    \item prestatiedrempels (bv. minimale detectienauwkeurigheid),
    \item vereisten rond generaliseerbaarheid naar variabele beeldcondities.
\end{itemize}

Door vooraf duidelijke criteria te bepalen, konden modellen op een objectieve manier worden beoordeeld.

\section*{Fase 3: Selectie van modellen en datasets}

In plaats van een brede long-list op te stellen, werd rechtstreeks overgegaan tot een gerichte selectie van modellen en datasets op basis van de bevindingen uit de literatuurstudie en de geformuleerde requirements.

De selectiecriteria omvatten onder andere:

\begin{itemize}
    \item beschikbaarheid als open-source implementatie,
    \item ondersteuning voor fine-tuning,
    \item toepasbaarheid in dense retailomgevingen,
\end{itemize}

Voor datasets werd gekeken naar representativiteit, annotatiekwaliteit en toepasbaarheid voor zowel training als evaluatie.

De output van deze fase was een afgebakende experimentele set-up bestaande uit geselecteerde modellen en bijhorende beelddata.

\section*{Fase 4: Ontwikkeling van een proof-of-concept}

Op basis van de geselecteerde modellen werd een proof-of-concept (PoC) ontwikkeld. Dit PoC had als doel de praktische inzetbaarheid van bestaande AI-modellen voor productdetectie in supermarktbeelden te demonstreren.

De ontwikkeling omvatte:

\begin{itemize}
    \item configuratie en eventuele fine-tuning van de modellen,
    \item implementatie van een detectiepipeline,
    \item verwerking van testbeelden,
    \item opslag van detectieresultaten voor verdere analyse.
\end{itemize}

De focus lag hierbij niet op het ontwikkelen van een nieuw model, maar op het experimenteel evalueren van bestaande architecturen binnen een concrete retailcontext.

\section*{Fase 5: Experimentele evaluatie}

De experimentele evaluatie werd uitgevoerd op beelden verzameld in realistische omstandigheden, zoals supermarktbeelden of gecontroleerde schapopstellingen (bv. voorraadkastscenario).

De testen werden opgezet om:

\begin{itemize}
    \item detectienauwkeurigheid te meten,
    \item prestaties te evalueren bij variërende rekdichtheid,
    \item robuustheid tegen veranderende camerahoek en afstand te analyseren,
    \item eventuele verschillen tussen modeltypes te identificeren.
\end{itemize}

De evaluatie gebeurde aan de hand van vooraf gedefinieerde metriek(en), zoals mAP en classificatienauwkeurigheid.

\section*{Fase 6: Analyse en conclusie}

In de laatste fase werden de experimentele resultaten statistisch en kwalitatief geanalyseerd. Hierbij werd nagegaan:

\begin{itemize}
    \item in welke mate de modellen voldeden aan de geformuleerde requirements,
    \item welke beeldfactoren significante invloed hadden op de prestaties,
    \item welke modelarchitectuur het meest geschikt bleek binnen de onderzochte context.
\end{itemize}

De analyse vormde de basis voor de uiteindelijke conclusie en aanbevelingen voor toekomstig onderzoek.

\chapter{\IfLanguageName{dutch}{Requirementsanalyse}{Requirements Analysis}}%
\label{ch:requirementsanalyse}

In dit hoofdstuk bespreken we modellen die we in de literatuurstudie al zijn tegengekomen en bespreken we de functionele vereisten, prestatievereisten en vereisten op het gebied van robuustheid en generaliseerbaarheid. Zo kunnen we de beste modellen voor onze toepassing selecteren.

\section{Functionele vereisten}

Op basis van de literatuurstudie kunnen we de volgende functionele vereisten opstellen. Het model moet:
\begin{itemize}
    \item producten kunnen detecteren in schapbeelden met hoge objectdensiteit.
    \item individuele producten kunnen lokaliseren door gebruik van bounding boxes.
    \item producten correct kunnen identificeren op SKU-niveau.
    \item kunnen omgaan met overlappende producten of scheve/beschadigde producten.
    \item visueel sterk gelijkende producten kunnen onderscheiden van elkaar.
    \item meerdere producten tegelijk kunnen verwerken binnen één beeld.
\end{itemize}


\section{Prestatievereisten}
\label{sec:prestatievereisten}

Op basis van prestatiedrempels kunnen we modellen objectief vergelijken en bekijken. We kiezen onze modellen op basis van volgende criteria:
\begin{itemize}
    \item Het model moet een mAP-score behalen op de testdataset van minstens 85\% en mAP50-95 van minstens 50\%.
    \item Het model moet een hoge recall behalen van minstens 80\%, aangezien het missen van producten vermeden moet worden.
    \item Het model moet een hoge precision behalen die minstens 85\% bedraagt, aangezien het model accuraat moet zijn over de herkenning van de producten.
    \item Voor classificatie moet een hoge recall@1 (Top-1 accuracy) worden nagestreefd die een drempelwaarde heeft van minstens 75\%.
    \item Recall@5 (Top-5 accuracy) wordt gebruikt als aanvullende metric voor visueel gelijkaardige producten.
\end{itemize}

\section{Robuustheid en generaliseerbaarheid}

Retailomgevingen variëren sterk, daarom moet het model robuust zijn tegen variaties in beeldcondities. Hiervoor kijken we naar:
\begin{itemize}
    \item variaties in belichting,
    \item verschillende camerahoeken en afstanden,
    \item variërende productdichtheid in schappen,
    \item en veranderingen in verpakking en productdesign.
\end{itemize}

De modellen moeten goed werken op beelden die lijken op de trainingsdata maar het is ook zeer belangrijk dat de modellen ook goed presteren op nieuwe of andere situaties.

\section{Vertaling van vereisten naar modelkeuze en pipeline}

De bovenstaande vereisten worden rechtstreeks vertaald naar de gekozen architectuur van de proof-of-concept. Geen enkel afzonderlijk model voldoet aan alle functionele en prestatievereisten tegelijk. Daarom wordt een combinatie van modellen gebruikt, waarbij elke component een specifieke taak uitvoert.

Voor objectdetectie kijken we naar YOLO-modellen. Deze modellen zijn geschikt voor real-time detectie en kunnen meerdere objecten tegelijk verwerken binnen één afbeelding. Het gebruik van bounding boxes sluit aan bij de functionele vereiste om individuele producten te lokaliseren. Er moet worden afgewogen tussen verschillende YOLO-varianten, waarbij de trade-off tussen nauwkeurigheid en snelheid in overweging wordt genomen.

Voor productidentificatie wordt niet gewerkt met een klassiek classificatiemodel, maar met een embedding-gebaseerde aanpak. In plaats van een vaste set klassen te voorspellen, wordt elk product voorgesteld als een vector in een hoge-dimensionale ruimte. Deze aanpak maakt het mogelijk om flexibel om te gaan met grote aantallen producten en ondersteunt het onderscheiden van visueel gelijkaardige items.

De embeddings worden gegenereerd met een combinatie van CLIP en Florence-2. CLIP levert een sterke algemene visuele representatie en presteert goed op globale kenmerken zoals kleur en merkidentiteit. Florence-2 vult dit aan met meer gedetailleerde visuele informatie, waardoor subtiele verschillen beter vastgelegd worden. Dit is noodzakelijk om te voldoen aan de vereiste om producten op SKU-niveau te onderscheiden.

Daarnaast wordt OCR toegepast via Florence-2 om tekstinformatie uit productverpakkingen te extraheren. Deze tekst wordt omgezet naar een embedding met behulp van de text encoder van CLIP. Hierdoor wordt expliciete productinformatie zoals merknamen en labels geïntegreerd in de representatie. Dit versterkt de identificatie, vooral bij producten die visueel sterk op elkaar lijken.

De combinatie van visuele en tekstuele embeddings wordt samengevoegd tot één vectorrepresentatie. Deze fusie zorgt ervoor dat verschillende informatiebronnen samen bijdragen aan de uiteindelijke vergelijking tussen producten.

Voor het zoeken naar overeenkomsten wordt gebruik gemaakt van FAISS. Deze component maakt het mogelijk om snel en efficiënt de meest gelijkaardige producten te vinden in een grote dataset. In plaats van een directe classificatie wordt gewerkt met een top-k benadering, waarbij meerdere mogelijke matches worden teruggegeven. Dit sluit aan bij het gebruik van Top-5 accuracy als aanvullende metric.

De gekozen pipeline ondersteunt verwerking van meerdere producten binnen één beeld. Elke detectie wordt afzonderlijk verwerkt, waardoor schaalbaarheid behouden blijft bij hoge objectdensiteit. Tegelijk blijft de kwaliteit afhankelijk van de nauwkeurigheid van de detectiestap en de representatiekracht van de embeddings.

Op deze manier wordt een systeem opgebouwd dat de functionele vereisten afdekt en tegelijkertijd rekening houdt met de beperkingen van individuele modellen.
\chapter{\IfLanguageName{dutch}{Selectie van modellen en datasets}{Selection of Models and Datasets}}%
\label{ch:selectie}


In dit hoofdstuk bespreken we de selectie van modellen en datasets die we zullen gebruiken in deze proef. We baseren onze selectie op de functionele vereisten, prestatievereisten en vereisten op het gebied van robuustheid en generaliseerbaarheid die we in hoofdstuk \ref{ch:requirementsanalyse} hebben opgesteld.

\section{Geselecteerde modellen}

Op basis van de literatuur en de gedefinieerde criteria wordt een selectie gemaakt van modellen uit verschillende categorieën.

\subsection{CNN-gebaseerde objectdetectie}

Als representatief model voor klassieke deep learning detectie wordt YOLOv8 gekozen \autocite{yolov8_ultralytics}. YOLO-modellen staan bekend om hun goede balans tussen snelheid en nauwkeurigheid en worden vaak gebruikt in real-time toepassingen \cite{ultralytics_docs}.

Volgens de literatuur worden CNN-modellen vaak gebruikt binnen voor retailtoepassingen, omdat ze zeer geschikt zijn voor het herkennen van producten \autocite{MELEK2024108452}. Ze zijn bovendien ook snel en betrouwbaar bij real-time systemen \cite{chadha2025vision}. 
Echter, CNN-modellen kunnen moeite hebben met het onderscheiden van sterk gelijkende producten, vooral in drukke beelden \autocite{PETTERSSON202410549}.

\subsection{Transformer-gebaseerde detectie}

Daarnaast wordt een transformer-gebaseerd model geselecteerd. Volgens de literatuur biedt dit type model voordelen bij complexe beelden waar vele objecten en overlappende producten in voorkomen \autocite{zhu2021}.

Modellen zoals RT-DETR maken gebruik van deze transformer-architectuur die de globale context beter begrijpt dan CNN-modellen, wat zeker ten goede komt bij beelden met hoge objectdensiteit zoals in een supermarkt \cite{carion2020endtoendobjectdetectiontransformers}.
Dit maakt het model bijzonder geschikt voor supermarktbeelden met hoge objectdensiteit.

Initieel werd er gekozen voor RT-DETR vanwege de sterke prestaties in objectdetectie, maar door gebrek aan hardwarecapaciteit werd uiteindelijk gekozen voor een lichtere variant, namelijk YOLOv11s \cite{ultralytics_docs}. 
Dit is een groot nadeel van transformer-gebaseerde modellen, ze vragen steeds meer rekenkracht en geheugen \cite{zhao2024detrsbeatyolosrealtime}, wat een uitdaging kan zijn bij het trainen en implementeren op consumentenhardware.
Moderne YOLO-architecturen maken steeds meer gebruik van aandachtmechanismen ("attention"). Dit helpt het model om zich te focussen op de meest relevante delen van een beeld en zijn geïnspireerd door transformer-architecturen \autocite{vaswani2023attentionneed}.
Hoewel YOLOv11s een hybride aanpak is, doordat het zowel CNN-gebaseerd als transformer-gebaseerd is, blijft het binnen de context van deze proef steeds relevant \cite{yolo11_ultralytics}.


\subsection{Multimodale benadering}

Om de beperkingen van visuele gelijkenis tussen producten te onderzoeken, wordt een vision-language model zoals CLIP bekeken.

Volgens \textcite{PETTERSSON202410549} kan het combineren van beeld en tekst helpen bij het onderscheiden van sterk gelijkende producten op SKU-niveau.

uitleg over waarom OWL-v2 gekozen werd van Google. Dit is ook te zwaar dus gebruik ik NanoOWL, een lichtere variant



lichtere varianten werden vaak gekozen door het gebrek aan hardware capaciteiten maar ook zodat de modellen uiteindelijk op consumentenhardware zou kunnen draaien zoals slimme camera's en handscanners.

\section{Geselecteerde datasets}

Er werd gekozen voor dataset SKU-110K \autocite{goldman2019dense} vanwege de hoge dichtheid aan objecten en de aanwezigheid van sterk gelijkende producten, wat aansluit bij de functionele vereisten van onze proef. Dit is een enorm grote dataset met 110.000 geannoteerde producten in 11.762 beelden, wat voldoende is voor training en evaluatie voor de gekozen modellen.

De SKU-110K dataset bevat train, test en validatie images, wat het mogelijk maakt om de modellen volledig te trainen maar ook te valideren en te testen. De beelden zijn afkomstig van echte winkels en simuleren dus zeer realistische omstandigheden.

% Om te testen hoe goed het model presteert op nieuwe beelden word er ook gebruik gemaakt van een tweede dataset, namelijk een zelf samengestelde dataset van beelden van Open Food Facts \autocite{openfoodfacts}. Hierdoor word getest of de modellen effectief kunnen herkennen of dit product in beeld is. 
\chapter{\IfLanguageName{dutch}{Ontwikkeling van een proof-of-concept}{Development of a Proof-of-Concept}}%
\label{ch:poc}
% Voeg hier je eigen hoofdstukken toe die de ``corpus'' van je bachelorproef
% vormen. De structuur en titels hangen af van je eigen onderzoek. Je kan bv.
% elke fase in je onderzoek in een apart hoofdstuk bespreken.

%\input{...}
%\input{...}
%...

%%=============================================================================
%% Conclusie
%%=============================================================================

\chapter{Conclusie}%
\label{ch:conclusie}

% TODO: Trek een duidelijke conclusie, in de vorm van een antwoord op de
% onderzoeksvra(a)g(en). Wat was jouw bijdrage aan het onderzoeksdomein en
% hoe biedt dit meerwaarde aan het vakgebied/doelgroep? 
% Reflecteer kritisch over het resultaat. In Engelse teksten wordt deze sectie
% ``Discussion'' genoemd. Had je deze uitkomst verwacht? Zijn er zaken die nog
% niet duidelijk zijn?
% Heeft het onderzoek geleid tot nieuwe vragen die uitnodigen tot verder 
%onderzoek?

Deze bachelorproef onderzocht in welke mate bestaande AI-modellen inzetbaar zijn voor automatische productherkenning in supermarktbeelden. De focus lag hierbij op de technische haalbaarheid van een reproduceerbare pipeline voor objectdetectie en similarity-based productidentificatie binnen retailomgevingen met hoge objectdichtheid.

Uit de literatuurstudie bleek dat automatische productherkenning in supermarkten een uitdagend probleem blijft. Producten staan vaak dicht bij elkaar, overlappen elkaar regelmatig en verschillen soms slechts in kleine details zoals kleuraccenten, tekst of verpakkingsvarianten. Daarnaast zorgen variaties in belichting, camerahoek en beeldkwaliteit voor bijkomende moeilijkheden. Hoewel moderne computer vision-modellen sterke prestaties tonen in algemene detectietaken, blijkt herkenning op SKU-niveau in realistische winkelomgevingen nog steeds complex.

Om deze problematiek te onderzoeken werd een proof-of-concept ontwikkeld bestaande uit meerdere componenten. Voor objectdetectie werd gebruik gemaakt van YOLO11s, terwijl productidentificatie gebeurde via embeddings gebaseerd op CLIP, OCR en Florence-2. De embeddings werden gecombineerd en vervolgens doorzocht via FAISS similarity search om mogelijke productmatches terug te geven.

De experimentele evaluatie toont aan dat fine-tuning een zeer grote impact heeft op detectieprestaties binnen retailomgevingen. Het generieke YOLO11s-model presteerde onvoldoende op supermarktbeelden met hoge objectdichtheid, maar na fine-tuning op de SKU-110K dataset steeg het aantal detecties aanzienlijk. Dit bevestigt dat modelspecialisatie noodzakelijk is voor toepassingen binnen complexe supermarktcontexten.

Daarnaast werd onderzocht welke combinatie van embeddings het meest geschikt is voor similarity-based productherkenning. Uit de experimenten bleek dat de combinatie van CLIP en OCR de meest consistente resultaten opleverde. Vooral tekstinformatie afkomstig van productverpakkingen bleek een belangrijke bijdrage te leveren aan correcte productmatching. De visuele encoder van Florence-2 bood daarentegen weinig meerwaarde binnen deze pipeline en werd uiteindelijk niet opgenomen in de finale configuratie.

Hoewel de top-5 resultaten vaak relevante producten bevatten, bleef de stabiliteit van top-1 voorspellingen beperkt. Producten met sterk gelijkende verpakkingen, zoals varianten van hetzelfde merk, bleken moeilijk betrouwbaar van elkaar te onderscheiden. Kleine verschillen in kleur, tekst of layout waren vaak onvoldoende om consistente embeddings te genereren. De vooropgestelde prestatiedrempels werden bovendien niet gehaald voor fine-grained productidentificatie, wat aantoont dat SKU-level herkenning in realistische retailomgevingen nog steeds een uitdagend probleem blijft. Dit benadrukt dat multimodale embeddings en similarity search reeds bruikbare resultaten opleveren, maar momenteel nog onvoldoende robuust zijn voor volledig betrouwbare identificatie op productniveau.

De resultaten tonen ook aan dat de kwaliteit van eerdere stappen in de pipeline een directe invloed heeft op het eindresultaat. Onnauwkeurige bounding boxes leiden tot slechte crops, wat vervolgens minder betrouwbare embeddings veroorzaakt. Daarnaast hadden OCR-fouten en inconsistente datasetafbeeldingen een negatieve impact op similarity matching. De prestaties van het volledige systeem worden dus bepaald door de interactie tussen detectie, embeddings, OCR en datasetkwaliteit.

Ondanks deze beperkingen toont deze bachelorproef aan dat automatische assortimentsinventarisatie technisch haalbaar is als onderzoeksrichting. Moderne computer vision-modellen kunnen reeds een groot deel van de producten in complexe supermarktbeelden detecteren en gedeeltelijk identificeren. Vooral de combinatie van objectdetectie met multimodale embeddings biedt interessante mogelijkheden voor toekomstige toepassingen binnen retailanalyse, voorraadbeheer en monitoring van gezonde of duurzame producten.

Toch blijven er nog verschillende beperkingen die verder onderzoek vereisen. De gebruikte datasets bevatten niet uitsluitend Vlaamse supermarktbeelden, waardoor de generaliseerbaarheid naar specifieke winkelcontexten beperkt blijft. Daarnaast bleef de evaluatie van retrievalprestaties relatief beperkt tot similarity scores en top-k matches. Hierdoor blijft het moeilijk om de retrievalkwaliteit volledig objectief te vergelijken met andere systemen. Toekomstig onderzoek kan uitgebreidere retrievalmetrics gebruiken, zoals Recall@k, Precision@k of Mean Reciprocal Rank, om productidentificatie nauwkeuriger te evalueren.

Verder onderzoek kan zich richten op krachtigere multimodale modellen, verbeterde OCR-methoden en grotere gespecialiseerde retaildatasets. Ook realtime verwerking op consumentenhardware blijft een belangrijk aandachtspunt wanneer dergelijke systemen praktisch inzetbaar moeten worden. Daarnaast kan onderzocht worden in welke mate gespecialiseerde retaildatasets uit Vlaamse supermarktcontexten de prestaties van SKU-level herkenning verder kunnen verbeteren.

Samenvattend kan gesteld worden dat bestaande AI-modellen reeds sterke mogelijkheden bieden voor automatische productdetectie in supermarktbeelden, maar dat betrouwbare productherkenning op SKU-niveau nog aanzienlijke uitdagingen bevat. Deze bachelorproef toont echter aan dat een combinatie van objectdetectie, embeddings en multimodale informatie een veelbelovende basis vormt voor verdere ontwikkeling binnen automatische retailanalyse.

%---------- Bijlagen -----------------------------------------------------------

\appendix

\chapter{Onderzoeksvoorstel}

Het onderwerp van deze bachelorproef is gebaseerd op een onderzoeksvoorstel dat vooraf werd beoordeeld door de promotor. Dat voorstel is opgenomen in deze bijlage.

%% TODO: 
%\section*{Samenvatting}

% Kopieer en plak hier de samenvatting (abstract) van je onderzoeksvoorstel.
In Vlaamse supermarkten is weinig transparantie beschikbaar over hoe duurzaam en gezond het werkelijk aangeboden assortiment is. Onderzoekers zeggen dat winkelomgevingen een bepalende invloed hebben op consumentengedrag en gezondheid, maar betrouwbare gegevens over het daadwerkelijke schapaanbod ontbreken grotendeels \autocite{wur2025, rug2024}. De overheid benadrukt de nood aan duidelijkere informatie rond duurzame voeding, maar stuit op het gebrek aan een schaalbare methode om producten in winkels objectief te registreren \autocite{rvo2025}.

In de praktijk gebeurt inventarisatie vandaag handmatig, wat tijdrovend, foutgevoelig en ongeschikt is voor grootschalige monitoring. Tegelijk evolueren technieken uit artificiële intelligentie, meer bepaald computer vision, snel. Modellen voor objectdetectie worden al getest in retail, maar de complexiteit van echte supermarktomgevingen, zoals hoge rekdichtheid, occlusies en sterk gelijkende verpakkingen, blijft een uitdaging \autocite{santra2019, fuchs2019}.

Dit onderzoek richt zich op het toepassen en evalueren van bestaande AI-modellen voor automatische productherkenning in Vlaamse supermarkten. De doelgroep omvat IT-professionals, data-engineers en onderzoekers die nood hebben aan een reproduceerbare, geautomatiseerde methode om supermarktassortimenten te analyseren. De centrale onderzoeksvraag luidt:

\begin{quote}
\textbf{In welke mate kunnen bestaande AI-modellen, al dan niet met bijkomende fine-tuning, producten in Vlaamse supermarktbeelden correct herkennen om zo transparantie te creëren over het voedingsaanbod?}
\end{quote}

% De doelstelling van deze bachelorproef is het ontwikkelen van een technisch onderbouwde vergelijkende studie, aangevuld met een proof-of-concept pipeline die productherkenning uitvoert en herkende items koppelt aan open databronnen. Een succesvol eindresultaat omvat een meetbaar inzicht in modelprestaties, een werkende prototype-implementatie en aanbevelingen voor verdere toepassing binnen onderzoek en beleid.
Om deze onderzoeksvraag te concretiseren, wordt het probleem verder opgesplitst in een aantal deelvragen binnen het probleemdomein:
\begin{itemize}
    \item Wat is het effect van gecontroleerde variaties in belichting (lux-niveau), beeldresolutie (px) en rekdichtheid (aantal producten per m²) op de mAP- en recall-scores van geselecteerde detectiemodellen?
    \item Welke bestaande open datasets of online databronnen bevatten geschikt beeldmateriaal van supermarktproducten voor training en evaluatie van productherkenningsmodellen?
\end{itemize}

Daarnaast worden deelvragen geformuleerd binnen het oplossingsdomein, die focussen op de technische haalbaarheid van automatische productherkenning:
\begin{itemize}
    \item Welke mAP@0.5:0.95-, recall- en precisionscores behalen geselecteerde AI-modellen bij detectie van producten in supermarktbeelden met hoge rekdichtheid?
    \item Hoe verschillen YOLO modellen in termen van mAP, recal en precision bij identieke supermarktbeelden?
    \item Wat is het effect van beperkte fine-tuning op mAP, recal en precision ten opzichte van het vooraf getrainde basismodel?
\end{itemize}

Om de centrale onderzoeksvraag objectief te beantwoorden, wordt gebruikgemaakt van gangbare evaluatiemetrieken uit objectdetectieonderzoek. Met mAP@0.5:0.95 krijg je een globaal beeld van correcte herkenningen, maar is onvoldoende in situaties met veel objecten per beeld. Precision en recall zijn daarom essentieel om respectievelijk fout-positieve detecties en gemiste producten te analyseren. 

De mean Average Precision (mAP) wordt gebruikt als kernmetriek, omdat deze zowel lokalisatie als classificatie over meerdere drempels combineert en toelaat modellen onderling te vergelijken in complexe omgevingen. Deze metriek sluit rechtstreeks aan bij de onderzoeksvraag, aangezien transparantie over het assortiment enkel zinvol is wanneer producten consistent en betrouwbaar worden herkend.

% Verwijzing naar het bestand met de inhoud van het onderzoeksvoorstel
%---------- Inleiding ---------------------------------------------------------

% TODO: Is dit voorstel gebaseerd op een paper van Research Methods die je
% vorig jaar hebt ingediend? Heb je daarbij eventueel samengewerkt met een
% andere student?
% Zo ja, haal dan de tekst hieronder uit commentaar en pas aan.

%\paragraph{Opmerking}

% Dit voorstel is gebaseerd op het onderzoeksvoorstel dat werd geschreven in het
% kader van het vak Research Methods dat ik (vorig/dit) academiejaar heb
% uitgewerkt (met medesturent VOORNAAM NAAM als mede-auteur).
% 

\section{Inleiding}%
\label{sec:inleiding}

In Vlaamse supermarkten is weinig transparantie beschikbaar over hoe duurzaam en gezond het werkelijk aangeboden assortiment is. Onderzoekers signaleren dat winkelomgevingen een bepalende invloed hebben op consumentengedrag en gezondheid, maar betrouwbare gegevens over het daadwerkelijke schapaanbod ontbreken grotendeels \autocite{wur2025, rug2024}. De overheid benadrukt de nood aan duidelijkere informatie rond duurzame voeding, maar stuit op het gebrek aan een schaalbare methode om producten in winkels objectief te registreren \autocite{rvo2025}.

In de praktijk gebeurt inventarisatie vandaag handmatig, wat tijdrovend, foutgevoelig en ongeschikt is voor grootschalige monitoring. Tegelijk evolueren technieken uit artificiële intelligentie, meer bepaald computer vision, snel. Modellen voor objectdetectie worden al getest in retail, maar de complexiteit van echte supermarktomgevingen, zoals hoge rekdichtheid, occlusies en sterk gelijkende verpakkingen, blijft een uitdaging \autocite{santra2019, fuchs2019}.

Dit onderzoek richt zich op het toepassen en evalueren van bestaande AI-modellen voor automatische productherkenning in Vlaamse supermarkten. De doelgroep omvat IT-professionals, data-engineers en onderzoekers die nood hebben aan een reproduceerbare, geautomatiseerde methode om supermarktassortimenten te analyseren. De centrale onderzoeksvraag luidt:

\begin{quote}
\textbf{In welke mate kunnen bestaande AI-modellen, al dan niet met bijkomende fine-tuning, producten in Vlaamse supermarktbeelden correct herkennen om zo transparantie te creëren over het voedingsaanbod?}
\end{quote}

% De doelstelling van deze bachelorproef is het ontwikkelen van een technisch onderbouwde vergelijkende studie, aangevuld met een proof-of-concept pipeline die productherkenning uitvoert en herkende items koppelt aan open databronnen. Een succesvol eindresultaat omvat een meetbaar inzicht in modelprestaties, een werkende prototype-implementatie en aanbevelingen voor verdere toepassing binnen onderzoek en beleid.
Om deze onderzoeksvraag te concretiseren, wordt het probleem verder opgesplitst in een aantal deelvragen binnen het probleemdomein:
\begin{itemize}
    \item Wat is het effect van gecontroleerde variaties in belichting (lux-niveau), beeldresolutie (px) en rekdichtheid (aantal producten per m²) op de mAP- en recall-scores van geselecteerde detectiemodellen?
    \item Welke bestaande open datasets of online databronnen bevatten geschikt beeldmateriaal van supermarktproducten voor training en evaluatie van productherkenningsmodellen?
\end{itemize}

Daarnaast worden deelvragen geformuleerd binnen het oplossingsdomein, die focussen op de technische haalbaarheid van automatische productherkenning:
\begin{itemize}
    \item Welke mAP@0.5:0.95-, recall- en SKU-level Top-1 accuracyscores behalen geselecteerde AI-modellen bij detectie van producten in supermarktbeelden met hoge rekdichtheid (>X producten per beeld)?
    \item Hoe verschillen YOLOv8 (gesloten-klassenset objectdetectie) en Grounding-DINO (open-vocabulary detectie) in termen van mAP, recall en detectiesnelheid (inference time per beeld) bij identieke supermarktbeelden?
    \item Wat is het effect van beperkte fine-tuning (≤10 epochs, enkel laatste detectielagen geüpdatet) op mAP- en SKU-level accuracy ten opzichte van het vooraf getrainde basismodel?
\end{itemize}

Om de centrale onderzoeksvraag objectief te beantwoorden, wordt gebruikgemaakt van gangbare evaluatiemetrieken uit objectdetectieonderzoek. Met mAP@0.5:0.95 krijg je een globaal beeld van correcte herkenningen, maar is onvoldoende in situaties met veel objecten per beeld. Precision en recall zijn daarom essentieel om respectievelijk fout-positieve detecties en gemiste producten te analyseren. 

De mean Average Precision (mAP) wordt gebruikt als kernmetriek, omdat deze zowel lokalisatie als classificatie over meerdere drempels combineert en toelaat modellen onderling te vergelijken in complexe omgevingen. Deze metriek sluit rechtstreeks aan bij de onderzoeksvraag, aangezien transparantie over het assortiment enkel zinvol is wanneer producten consistent en betrouwbaar worden herkend.


%---------- Stand van zaken ---------------------------------------------------

\section{Literatuurstudie}%
\label{sec:literatuurstudie}

Automatische productherkenning in retailomgevingen is een actief onderzoeksdomein binnen computer vision. Hoewel commerciële oplossingen bestaan, blijkt uit onderzoek dat realistische supermarktcontexten bijzonder uitdagend zijn vanwege variabele lichtomstandigheden, occlusies, overlappingen en de sterke visuele gelijkenis tussen verpakkingen \autocite{santra2019}. Retailonderzoek draait vaak rond objectdetectie met convolutionele neurale netwerken of transformer-gebaseerde modellen zoals YOLO, EfficientDet en Vision Transformers.

\textcite{fuchs2019} tonen aan dat deep-learning-modellen goed presteren in gecontroleerde omstandigheden, maar dat prestaties dalen wanneer rekdichtheid en variabele camerastandpunten toenemen. Ook generaliseerbaarheid naar nieuwe productvarianten blijft een probleem. Recente modellen zoals YOLOv8, Grounding-DINO en CLIP-gebaseerde pipelines combineren detectie en multimodale embedding-technieken om robuustere herkenning mogelijk te maken, maar hun prestaties in de Belgische of Vlaamse supermarktcontext zijn nog niet onderzocht.

Daarnaast is de koppeling met externe databronnen een belangrijk aspect. OpenFoodFacts is internationaal gebruikt in diverse voedingsstudies, maar de volledigheid en correctheid van bepaalde metadata, zoals eco-scores en detailingrediënten, blijkt wisselend, wat uitdagingen creëert voor automatische productkoppeling \autocite{wur2025}.

Uit de literatuur blijkt dat er een duidelijke lacune is: weinig studies richten zich op lokale supermarktcontexten, en er is een gebrek aan evaluaties die modellen vergelijken onder realistische omstandigheden. Dit onderzoek sluit hierop aan door een experimentele studie uit te voeren binnen Vlaamse supermarkten en de prestaties van meerdere AI-modellen systematisch te toetsen.

% Voor literatuurverwijzingen zijn er twee belangrijke commando's:
% \autocite{KEY} => (Auteur, jaartal) Gebruik dit als de naam van de auteur
%   geen onderdeel is van de zin.
% \textcite{KEY} => Auteur (jaartal)  Gebruik dit als de auteursnaam wel een
%   functie heeft in de zin (bv. ``Uit onderzoek door Doll & Hill (1954) bleek
%   ...'')

%---------- Methodologie ------------------------------------------------------
\section{Methodologie}%
\label{sec:methodologie}

Het onderzoek volgt een experimenteel-opgezette vergelijkende studie. Elke fase van het stappenplan levert observaties of meetresultaten op die rechtstreeks terugkoppelen naar de centrale onderzoeksvraag, namelijk de mate waarin bestaande AI-modellen bruikbaar zijn voor het creëren van transparantie over het voedingsaanbod in Vlaamse supermarkten. De methodologie bestaat uit vier fases:

\subsection*{1. Verzamelen en voorbereiden van beelddata}
Tijdens deze fase wordt een systematische verkenning uitgevoerd van bestaande open datasets en publiek toegankelijke online databronnen met productbeelden. Deze bronnen worden geëvalueerd op resolutie, labelkwaliteit en beschikbaarheid van metadata. De geschiktheid van deze datasets voor de Vlaamse supermarktcontext wordt beoordeeld op basis van producttypes, verpakkingsvormen en schapstructuren.
Alle data worden geanonimiseerd en verwerkt volgens ethische richtlijnen.

\subsection*{2. Selectie en configuratie van AI-modellen}
De volgende modellen worden geselecteerd op basis van hun geschiktheid voor retaildetectie:
\begin{itemize}
    \item YOLOv8 (objectdetectie, hoge snelheid),
    \item Grounding-DINO (open-vocab detection),
    \item CLIP-gebaseerde pipelines (multimodale classificatie).
\end{itemize}
De selectie van zowel een klassiek objectdetectiemodel als open-vocab en multimodale modellen laat toe om verschillen in detectiegedrag te observeren. Door deze modellen onder identieke omstandigheden toe te passen, kan onderzocht worden in welke mate de modelkeuze invloed heeft op de betrouwbaarheid van automatische productherkenning in supermarktcontexten.

\subsection*{3. Experimentele evaluatie}
Voor elk model worden prestaties gemeten onder variabele omstandigheden. De detectieresultaten van de modellen worden onder identieke datasets en instellingen vergeleken om verschillen in robuustheid en foutpatronen bloot te leggen. Voor en na fine-tuning worden de prestaties vergeleken om het effect van beperkte lokale training op detectienauwkeurigheid te kwantificeren. De volgende metrieken worden systematisch gemeten:
\begin{itemize}
    \item mAP@0.5:0.95,
    \item Recall,
    \item SKU-level Top-1 accuracyscores.
\end{itemize}
De gekozen evaluatiemetrieken zijn afgestemd op de centrale onderzoeksvraag. mAP@0.5:0.95 wordt enkel indicatief gebruikt en niet als beslissende metriek, aangezien deze onvoldoende robuust is in beelden met meerdere objecten. SKU-level Top-1 accuracyscores en recall worden daarom afzonderlijk geanalyseerd om respectievelijk fout-positieve detecties en gemiste producten in kaart te brengen, wat essentieel is voor transparantie over het werkelijke assortiment.

De mean Average Precision (mAP) fungeert als kernmetriek, aangezien deze zowel lokalisatie als classificatie over meerdere drempels combineert en toelaat modellen onderling te vergelijken. Een hogere mAP-score wijst op consistente en reproduceerbare detectie, wat een noodzakelijke voorwaarde is om automatische inventarisatie bruikbaar te maken voor onderzoek en beleid.
Modellen met lage SKU-level Top-1 accuracyscores of recall worden als onvoldoende beschouwd om het volledige assortiment betrouwbaar in kaart te brengen, aangezien zowel over- als onderdetectie de transparantie van het aanbod vertekenen.

De prestaties worden niet enkel globaal geëvalueerd, maar ook uitgesplitst per beeldfactor (licht, resolutie, camerahoek en rekdichtheid) om hun invloed op de herkenningsnauwkeurigheid te kwantificeren.

\subsection*{4. Koppeling van herkende producten aan open databronnen}
Deze fase beantwoordt de deelvraag in welke mate detectieresultaten praktisch bruikbaar zijn voor verdere analyse. Voor elk herkend product wordt nagegaan of voldoende informatie beschikbaar is om een overeenkomst te vinden met een item in OpenFoodFacts. Het aandeel succesvolle koppelingen en de aard van mislukte pogingen worden gerapporteerd als indicatie van de toepasbaarheid van automatische productherkenning in combinatie met open databronnen.
Hierbij wordt gebruikgemaakt van barcodes, tekstextractie (OCR) en visuele embeddings.

\subsection*{Planning en deliverables}
\begin{itemize}
    \item Week 1–3: literatuurstudie, dataverzameling.
    \item Week 4–7: modelselectie, opzetten testomgeving, eerste metingen.
    \item Week 8–10: fine-tuning en uitgebreide evaluatie.
    \item Week 11–12: implementatie koppeling met open databronnen.
    \item Week 13–14: analyse, rapportage, aanbevelingen.
\end{itemize}

Door modelprestaties, gevoeligheid voor beeldfactoren, het effect van fine-tuning en de praktische koppelbaarheid aan open databronnen gezamenlijk te analyseren, laat deze methodologie toe om objectief te bepalen in welke mate bestaande AI-modellen inzetbaar zijn voor transparante en schaalbare inventarisatie van het supermarktassortiment.
%---------- Verwachte resultaten ----------------------------------------------
\section{Verwacht resultaat, conclusie}%
\label{sec:verwachte_resultaten}

De resultaten zijn afhankelijk van de beschikbaarheid en kwaliteit van beelddata en kunnen niet zonder meer worden veralgemeend naar alle supermarktketens.

De verwachting is dat de performantie van de geselecteerde modellen sterk varieert afhankelijk van beeldomstandigheden. YOLOv8 wordt vermoedelijk het snelst met stabiele baseline-prestaties, terwijl Grounding-DINO mogelijk betere herkenning biedt bij sterk variërende verpakkingen maar trager werkt. Fine-tuning op Vlaamse beelden zal naar verwachting een merkbare verbetering opleveren in mAP-scores, vooral voor sterk gelijkende producten.

Een voorbeeld van een te verwachten grafiek is een vergelijkende staafdiagram met op de X-as de verschillende modellen (YOLOv8, CLIP-gebaseerde pipelines, Grounding-DINO) en op de Y-as de mAP-score per conditie (normaal licht, lage resolutie, schuine hoek). Deze visualiseert welke modellen het meest robuust blijven onder realistische omstandigheden.

De beoogde meerwaarde voor de doelgroep bestaat uit:
\begin{itemize}
    \item een objectieve evaluatie van AI-modellen specifiek in Vlaamse supermarktcontexten,
    \item een reproduceerbare testopzet die door andere onderzoekers hergebruikt kan worden,
    \item een proof-of-concept dat beleidsmakers en onderzoekers toelaat om het supermarktassortiment automatisch in kaart te brengen.
\end{itemize}

Zelfs wanneer bepaalde modellen onvoldoende presteren, biedt dit onderzoek waardevolle inzichten in hun beperkingen en vereisten voor verdere ontwikkeling. De bachelorproef levert daarmee directe praktische en wetenschappelijke waarde op voor voedingsonderzoek, data-engineering en retailinnovatie.




%%---------- Andere bijlagen --------------------------------------------------
% TODO: Voeg hier eventuele andere bijlagen toe. Bv. als je deze BP voor de
% tweede keer indient, een overzicht van de verbeteringen t.o.v. het origineel.
%\input{...}

%%---------- Backmatter, referentielijst ---------------------------------------

\backmatter{}

\setlength\bibitemsep{2pt} %% Add Some space between the bibliograpy entries
\printbibliography[heading=bibintoc]

\end{document}
